% Preamble templated from Mihir-Divyansh/Course-Setup
%iffalse
\let\negmedspace\undefined
\let\negthickspace\undefined
\documentclass[journal,12pt,onecolumn]{IEEEtran}
\usepackage{cite}
\usepackage{amsmath,amssymb,amsfonts,amsthm}
\usepackage{algorithmic}
\usepackage{graphicx}
\usepackage{textcomp}
\usepackage{xcolor}
\usepackage{txfonts}
\usepackage{listings}
\usepackage{enumitem}
\usepackage{mathtools}
\usepackage{gensymb}
\usepackage{comment}
\usepackage[breaklinks=true]{hyperref}
\usepackage{tkz-euclide}
\usepackage{listings}
\usepackage{gvv}
%\def\inputGnumericTable{}
\usepackage[latin1]{inputenc}
\usepackage{color}
\usepackage{array}
\usepackage{longtable}
\usepackage{calc}
\usepackage{multirow}
\usepackage{hhline}
\usepackage{ifthen}
\usepackage{lscape}
\usepackage{tabularx}
\usepackage{array}
\usepackage{float}
\usepackage{caption}
\usepackage{multicol}

\newtheorem{theorem}{Theorem}[section]
\newtheorem{problem}{Problem}
\newtheorem{proposition}{Proposition}[section]
\newtheorem{lemma}{Lemma}[section]
\newtheorem{corollary}[theorem]{Corollary}
\newtheorem{example}{Example}[section]
\newtheorem{definition}[problem]{Definition}
\newcommand{\BEQA}{\begin{eqnarray}}
\newcommand{\EEQA}{\end{eqnarray}}
%\newcommand{\define}{\stackrel{\triangle}{=}}
\theoremstyle{remark}
%\newtheorem{rem}{Remark}

% Marks the beginning of the document
\begin{document}
\bibliographystyle{IEEEtran}
\vspace{3cm}

\title{Signal Processing: Fourier Transform Time Delay}
\author{EE25BTECH11055 - Subhodeep Chakraborty}
\maketitle
\hrulefill
\bigskip

\renewcommand{\thefigure}{\theenumi}
\renewcommand{\thetable}{\theenumi}

\textbf{Question:}\par
Let the signal $x(t)$ have the Fourier transform $X(\omega)$. Consider the signal $y(t) = x(t-t_a)$ where $t_a$ is an arbitrary delay. The magnitude of the Fourier transform of $y(t)$ is:
\begin{enumerate}[label=\alph*)]
    \item $|X(\omega) \cdot [\omega]|$
    \item $|X(\omega)| \cdot \omega$
    \item $\omega^2 \cdot |X(\omega)|$
    \item $|X(\omega) \cdot e^{-j\omega t_a}|$
\end{enumerate}

\textbf{Solution:}\par

Let $x(t)$ be mapped to the frequency domain via the \textbf{Continuous Fourier Transform Operator} $\mathcal{F}$, such that $X(\omega) = \mathcal{F}\{x(t)\}$.

We are given a time-shifted signal $y(t) = x(t - t_a)$. According to the \textbf{Time-Shifting Theorem of the Fourier Transform}, a delay in the time domain corresponds to a phase rotation in the frequency domain.

To formalize this using linear algebra, we can represent the continuous frequency spectrum at discrete frequency bins $\omega_k$ as a column vector $\mathbf{X}$. The time-shift operation acts as a \textbf{Linear Transformation} mapping input spectrum $\mathbf{X}$ to output spectrum $\mathbf{Y}$. Because the complex exponentials are eigenfunctions of Linear Time-Invariant (LTI) systems, this transformation is represented by a diagonal \textbf{Phase-Shift Matrix} $\mathbf{\Phi}$:

\begin{align}
    \mathbf{Y} &= \mathbf{\Phi} \mathbf{X}
\end{align}

Expanding this matrix equation yields:
\begin{align}
    \begin{bmatrix} Y(\omega_1) \\ Y(\omega_2) \\ \vdots \\ Y(\omega_k) \end{bmatrix}
    &=
    \begin{bmatrix}
    e^{-j\omega_1 t_a} & 0 & \dots & 0 \\
    0 & e^{-j\omega_2 t_a} & \dots & 0 \\
    \vdots & \vdots & \ddots & \vdots \\
    0 & 0 & \dots & e^{-j\omega_k t_a}
    \end{bmatrix}
    \begin{bmatrix} X(\omega_1) \\ X(\omega_2) \\ \vdots \\ X(\omega_k) \end{bmatrix}
\end{align}

To find the magnitude of the resulting spectrum $|\mathbf{Y}|$, we take the element-wise absolute value. For any arbitrary frequency component $\omega$, the operation simplifies to:
\begin{align}
    |Y(\omega)| &= |X(\omega) \cdot e^{-j\omega t_a}| \\
    |Y(\omega)| &= |X(\omega)| \cdot |e^{-j\omega t_a}| \label{eq:magnitude_split}
\end{align}

To evaluate the magnitude of the complex exponential operator $|e^{-j\omega t_a}|$, we apply \textbf{Euler's Formula} ($e^{j\theta} = \cos\theta + j\sin\theta$) to express it in rectangular coordinates:
\begin{align}
    |e^{-j\omega t_a}| &= |\cos(-\omega t_a) + j\sin(-\omega t_a)|
\end{align}

Applying the \textbf{Pythagorean Trigonometric Identity} ($\cos^2\theta + \sin^2\theta = 1$) to find the modulus of this complex number:
\begin{align}
    |e^{-j\omega t_a}| &= \sqrt{\cos^2(-\omega t_a) + \sin^2(-\omega t_a)} = \sqrt{1} = 1
\end{align}

Substituting this unity magnitude back into Equation \eqref{eq:magnitude_split} gives:
\begin{align}
    |Y(\omega)| &= |X(\omega)| \cdot 1 \\
    |Y(\omega)| &= |X(\omega)|
\end{align}

Looking at the given options, option (d) represents the direct mathematical equivalence before simplification:
\begin{align}
    |X(\omega) \cdot e^{-j\omega t_a}| &= |X(\omega)| \cdot |e^{-j\omega t_a}| = |X(\omega)|
\end{align}

Therefore, the magnitude of the signal's frequency spectrum is entirely invariant to linear time delays. The correct choice is \textbf{(d)}.

\end{document}
